\chapter{Dysgraphia}
\index{Dysgraphia}

\section*{Definition and Overview}
Dysgraphia is a learning difficulty that affects a person's ability to write. It involves persistent difficulty with handwriting, spelling, and the organisation of written expression that is not explained by intelligence, vision, or general educational opportunity. Writing may be slow, effortful, messy, or poorly organised, and the task of physically producing letters can demand so much attention that the person cannot think about what they want to say. Dysgraphia is a lifelong condition, but with early recognition, appropriate teaching, and the use of technology, most people learn to manage writing tasks successfully.

\section*{Epidemiology}
Dysgraphia is estimated to affect a small but significant proportion of school-aged children, with figures often quoted at around 5 to 10 per cent of students. It is more common in boys and frequently coexists with other learning difficulties, particularly dyslexia and attention deficit hyperactivity disorder. Like other specific learning difficulties, it is often under-recognised, and children may be labelled as careless or unmotivated rather than given appropriate support.

\section*{Aetiology and Pathophysiology}
The exact cause of dysgraphia is not fully understood, but it involves difficulties in the brain processes underlying writing.
\begin{itemize}
    \item \textbf{Motor planning:} difficulty planning and executing the fine motor movements needed to form letters
    \item \textbf{Visual-spatial processing:} trouble with the spatial organisation of letters and words on the page
    \item \textbf{Working memory:} difficulty holding words and ideas in mind while coordinating the physical act of writing
    \item \textbf{Language processing:} difficulties with spelling and the translation of thoughts into written language
    \item \textbf{Genetic factors:} a family tendency, often alongside other learning difficulties
    \item \textbf{Coexisting conditions:} dyslexia, ADHD, and autism frequently occur together with dysgraphia
    \item \textbf{Brain injury:} acquired dysgraphia can follow stroke or head injury affecting the writing areas of the brain
\end{itemize}

\section*{Clinical Features}
The signs of dysgraphia appear when a child begins formal writing tasks and persist, in some form, throughout life.
\begin{itemize}
    \item Illegible or inconsistent handwriting, with poorly formed letters of uneven size
    \item Slow, laboured writing that is tiring and painful in the hand
    \item Unusual pencil grip, wrist position, or body posture
    \item Difficulty spacing words and staying on the line
    \item Frequent crossing out and erasing
    \item Poor spelling and difficulty learning letter sounds and patterns
    \item Writing far shorter pieces than the child can express verbally
    \item Difficulty organising ideas in written form, despite good verbal skills
    \item Avoidance of writing tasks and frustration with schoolwork
\end{itemize}

\section*{Differential Diagnosis}
Writing difficulties have a range of possible causes that should be explored.
\begin{itemize}
    \item Developmental coordination disorder, which affects general motor skills
    \item Dyslexia, which primarily affects reading and spelling
    \item Attention deficit hyperactivity disorder, affecting task focus and organisation
    \item Intellectual disability or global learning difficulties
    \item Vision problems or hand-eye coordination issues
    \item Emotional or motivational factors, which may mask or mimic dysgraphia
\end{itemize}

\section*{When to Seek Medical Care}
Parents and teachers should seek assessment when a child's writing is significantly behind expectations, causes distress, or interferes with learning. A GP can coordinate referral to a paediatrician, psychologist, or occupational therapist for assessment.

\textbf{Contact your local emergency services for:}
\begin{itemize}
    \item Sudden loss of writing ability or other new language or motor problems in an adult, which may indicate a stroke
    \item Sudden weakness, confusion, or difficulty speaking
\end{itemize}

\section*{Diagnosis and Investigations}
Diagnosis of dysgraphia involves a multidisciplinary assessment.
\begin{itemize}
    \item \textbf{Educational history:} samples of writing, reports of effort, and performance across subjects
    \item \textbf{Occupational therapy assessment:} evaluation of fine motor skills, pencil grip, and handwriting speed and legibility
    \item \textbf{Psychological assessment:} cognitive testing to identify the pattern of strengths and difficulties
    \item \textbf{Standardised writing tests:} measures of handwriting, spelling, and written expression
    \item \textbf{Screening for coexisting conditions:} assessment for dyslexia and ADHD
\end{itemize}

\section*{Management}
Management is educational, therapeutic, and practical.
\begin{itemize}
    \item \textbf{Occupational therapy:} exercises and strategies to improve fine motor control, grip, and letter formation
    \item \textbf{Handwriting programmes:} structured teaching of letter formation and automatic writing
    \item \textbf{Keyboard skills:} teaching touch typing, which removes the motor burden of handwriting
    \item \textbf{Assistive technology:} word processors, speech-to-text software, and writing apps
    \item \textbf{School accommodations:} extra time, reduced copying demands, and alternatives to handwritten assessment
    \item \textbf{Support for written expression:} teaching planning, drafting, and organising strategies
    \item \textbf{Emotional support:} building confidence and reducing frustration and avoidance
\end{itemize}

\section*{Complications}
\begin{itemize}
    \item Low self-esteem, frustration, and avoidance of schoolwork
    \item Falling behind academically despite strong verbal ability
    \item Misunderstanding and criticism from teachers and family
    \item Anxiety, particularly around written examinations
    \item Hand pain and fatigue from effortful writing
    \item Reduced career options if support is not provided
\end{itemize}

\section*{Prognosis and Outcomes}
Dysgraphia is a lifelong condition, but outcomes are generally good with early identification and support. Many children learn to write legibly and fluently enough for daily needs, and adults with dysgraphia manage written tasks effectively through keyboards and assistive technology. The key is recognising the difficulty early, adjusting expectations, and focusing on strategies rather than struggling with unmodified handwriting demands.

\section*{Prevention and Self-Care}
\begin{itemize}
    \item Seek assessment early rather than waiting for years of struggle
    \item Encourage writing through fun, low-pressure activities
    \item Teach keyboard skills from primary school
    \item Use speech-to-text and other tools without stigma
    \item Focus on content and ideas rather than neatness in the early stages
    \item Advocate for accommodations in school and examinations
    \item Praise effort and progress, not just final appearance
\end{itemize}

\section*{Sources and Further Reading}
The following reputable sources provide further information for healthcare students and patients seeking reliable, current advice about dysgraphia.
\begin{itemize}
    \item SPELD NSW and other state associations. \textit{Dysgraphia information and support.} https://www.speldnsw.org.au/
    \item Australian Occupational Therapy Association. \textit{Occupational therapy for handwriting difficulties.} https://www.otaus.com.au/
    \item Understood.org. \textit{Dysgraphia: what it is and how to help.} https://www.understood.org/
    \item Mayo Clinic. \textit{Learning disorders: symptoms and causes.} https://www.mayoclinic.org/diseases-conditions/learning-disorders/
    \item NHS. \textit{Learning disabilities in children.} https://www.nhs.uk/conditions/learning-disabilities/
    \item International Dyslexia Association. \textit{Dysgraphia resources.} https://dyslexiaida.org/
\end{itemize}
